\documentclass[a4paper]{article}

\usepackage[english]{babel}
%\usepackage[utf8x]{inputenc}
\usepackage{amsmath}
\usepackage{graphicx}
\usepackage[colorinlistoftodos]{todonotes}
\usepackage[all]{xy}
\usepackage{dsfont}
\usepackage{amssymb}
\usepackage{bbold}
\usepackage{amsthm}

\usepackage{tikz}
\usepackage{verbatim}
\usetikzlibrary{matrix}
\usetikzlibrary{cd}

\usepackage[autostyle=true]{csquotes}
\usepackage{pdfpages}
\usepackage{hyperref}
\usepackage{fancyhdr}



\newtheorem{theorem}{Theorem}[section]
\newtheorem{proposition}{Proposition}[theorem]
\newtheorem{lemma}{Lemma}[theorem]
\newtheorem{corollary}{Corollary}[theorem]
\newtheorem*{definition}{Definition}


\newcommand{\norm}[1]{\left\lVert#1\right\rVert}

\newcommand{\Rn}{\mathbb{R}^n}


\newcommand{\too}[1]{\overset{#1}{\longrightarrow}}
\newcommand{\hookar}[1]{\overset{#1}{\hookrightarrow}}
\newcommand{\colim}[1]{ \underset{#1}{\texttt{colim}} \text{ } }
\newcommand{\fset}[1]{ \{ #1_0 , \dots , #1_n \} }
\newcommand{\ddt}{\frac{d}{dt}}


\newcommand{\N}{\mathbb{N}}
\newcommand{\Z}{\mathbb{Z}}







\fancyhead[L]{\includegraphics[scale=0.5]{UFT.png}}

\begin{document}

\begin{center}
    \Large{\textbf{Soluciones}}
\end{center}

\setlength{\parindent}{0cm}


\textbf{Ejercicio 1:} Encuentre los máximos locales, mínimos locales y puntos silla de la función

\begin{align*}
    f(x,y) &=  \frac{x^2y^2-8x+y}{xy}\\
\end{align*}

\textbf{Solución: } Notamos que

\begin{align*}
    \frac{x^2y^2-8x+y}{xy} &= \frac{x^2y^2}{xy}-\frac{8x}{xy}+\frac{y}{xy} \\
    &= xy-\frac{8}{y}+\frac{1}{x}
\end{align*}

Calculamos el gradiente


\begin{align*}
    \nabla(f) = \left( y-\frac{1}{x^2}, x+\frac{8}{y^2} \right)
\end{align*}

Luego $\nabla(f)= (0,0)$, implica

\begin{align*}
     y-\frac{1}{x^2} &= 0 \\
     x+\frac{8}{y^2} &= 0
\end{align*}

De la primera ecuación, obtenemos que $\displaystyle y = \frac{1}{x^2}$, reemplazando en la segunda


\begin{align*}
    x+\frac{8}{\left(\frac{1}{x^2}\right)^2} &= 0 \\
    x+8x^4 &= 0 \\
    x(1+8x^3) &= 0
\end{align*}

De donde se deduce que $x=0$ o bien $1+8x^3=0$, pero $x=0$ es imposible pues $y=1/x^2$, luego la única opción posible es $1+8x^3=0$, de donde se concluye que

\begin{align*}
    x^3 &= -\frac{1}{8} \\
    x & = -\frac{1}{2}
\end{align*}

De esta forma 

\begin{align*}
    y=4
\end{align*}

Por lo tanto el punto crítico es $(-1/2, 4)$. Calculamos el Hessiano

\begin{align*}
    H(x,y) = \begin{vmatrix}
        \frac{2}{x^3} & 1 \\
        1  &  -\frac{16}{y^3}
    \end{vmatrix}
\end{align*}

Evaluamos en $(-1/2, 4)$

\begin{align*}
    H(-1/2,4) = \begin{vmatrix}
        -16 & 1 \\
        1  &  -\frac{1}{4}
    \end{vmatrix} = 4-1 =3
\end{align*}

Notamos que $-16 < 0$ mientras que $ H(-1/2,4)  > 0$, concluimos que $(-1/2,4)$ es un máximo local







\newpage



\textbf{Ejercicio 2:} Utilice los multiplicadores de Lagrange para encontrar los valores extremos de la función $f(x,y)=e^{xy}$ , con la restricción $x^2+y^2=16$ \\

\textbf{Solución: }Consideramos $g(x,y)= x^2+y^2-16$, tenemos que resolver $\nabla f = \lambda \nabla g$, es decir

\begin{align*}
    (ye^{xy},xe^{xy}) &= \lambda(2x,2y)
\end{align*}

Luego

\begin{align*}
    ye^{xy} &= 2\lambda x \\
    xe^{xy} &= 2\lambda y \\
\end{align*}

Dividimos las ecuaciones


\begin{align*}
   \frac{y}{x} &= \frac{x}{y}
\end{align*}

Luego $x^2=y^2$, reemplazamos en $x^2+y^2=16$ y obtenemos $2x^2=16$, luego $x^2=8$ y por lo tanto

\begin{align*}
    x=\pm 2\sqrt{2}
\end{align*}

De la misma forma para cada valor de $x$, obtenemos $y= \pm\sqrt{2}$. Por lo tanto tenemos cuatro extremos $(2\sqrt{2},2\sqrt{2})$, $(2\sqrt{2}, - 2\sqrt{2})$, $(-2\sqrt{2},2\sqrt{2})$ y $(-2\sqrt{2},-2\sqrt{2})$





\newpage

\textbf{Ejercicio 3:} \\

a) Sea $R$ la región triangular del plano cuyos vértices son los puntos de coordenadas $(0,0)$, $(1,1)$ y $(0,1)$. Exprese el área de la región $R$ usando integrales dobles, en los órdenes $dydx$ y $dxdy$ \\

b)  Calcule el área de la superficie $f(x,y)= 3x+y^2$, que se encuentra encima de la región $R$. \\

\textbf{Solución: } a) Las integrales quedan de la forma

\begin{align*}
    \int_0^1\int_x^1 dydx 
\end{align*}

\begin{align*}
    \int_0^1\int_0^y dxdy 
\end{align*}

b)  Tenemos que $f_x=3$ y $f_y=2y$, luego queremos calcular

\begin{align*}
    A = \iint_R \sqrt{10+4y^2}dA
\end{align*}

Para ello usamos la integral de la forma $dxdy$


\begin{align*}
    A &= \int_0^1\int_0^y   \sqrt{10+4y^2}dxdy \\
    &= \int_0^1  y\sqrt{10+4y^2}dy \\
\end{align*}

Ahora usamos la sustitución $u=10+4y^2$ luego $dy=\frac{du}{8}$, luego tenemos


\begin{align*}
    A &= \frac{1}{8}\int_10^14 \sqrt{u}du \\
    &= \frac{1}{8}\int_10^14 \sqrt{u}du \\
    &= \frac{1}{8} \frac{2x^{3/2}}{3} \bigg|_{10}^{14} \\
    &= \frac{14\sqrt{14}-10\sqrt{10}}{12}
\end{align*}


\newpage

\textbf{Ejercicio 4:} Calcular

\begin{align*}
    \iint_V (2zx^2+2zy^2)dxdydz
\end{align*}

Donde $V$ es el volumen del sólido exterior a la hoja superior del cono $z=\sqrt{x^2+y^2}$ e interior al cilindro $x^2+y^2=1$, con $z \geq 0$ \\

\textbf{Solución: } Observando el dibujo y pasando a cilíndircas obtenemos que

\begin{align*}
    V &= \int_0^{2\pi}\int_0^1\int_0^{r} 2z(r^3\cos^2(\theta)+r^2\sin^2(\theta))dzdrd\theta \\
    &= \int_0^{2\pi}\int_0^1 r^3 \int_0^{r} 2zdzdrd\theta \\
    &= \int_0^{2\pi}\int_0^1 r^5drd\theta \\
    &= \int_0^{2\pi} \frac{1}{6}d\theta \\
    &= \frac{1}{6}2\pi \\
    &= \frac{\pi}{3}
\end{align*}




\end{document}